\documentclass[10pt,a4paper]{amsart}
\usepackage[margin=19mm]{geometry}
\usepackage{amsmath,amssymb,mathtools}
\usepackage[scr=rsfs,cal=euler]{mathalfa}
\usepackage{xfrac}
\usepackage[unicode,hidelinks,bookmarks=false]{hyperref}
\newcommand{\ie}{i.e.\ }
\newcommand{\cf}{cf.\ }
\newcommand{\resp}{respectively\ }
\newcommand{\Subsection}{\subsection}
\providecommand{\nobreakdash}{\nobreak\hbox{-}}
\setlength{\emergencystretch}{2em}
\multlinegap0pt
\usepackage{bm}
\usepackage{bbm}
\usepackage{dsfont}
\usepackage{xspace}
\usepackage{cancel}
\usepackage{mathtools}
\usepackage{amsthm}
\makeatletter
\@ifundefined{thm}{\newtheorem{thm}{Thm}}{}
\@ifundefined{proposition}{\newtheorem{proposition}[thm]{Proposition}}{}
\makeatother
\DeclareMathOperator*{\JI}{JI}
\DeclareMathOperator*{\LM}{LM}
\DeclareMathOperator*{\MI}{MI}
\title[Left modularity and extremality for (some) infinite lattices]{Left modularity and extremality for (some) infinite lattices}
\author{Sota Asai, Osamu Iyama, Kaveh Mousavand, Charles Paquette}
\date{Source: arXiv:2604.20947v1 (CC BY 4.0)}
\begin{document}
\maketitle

\noindent\emph{Source and licence.} Left modularity and extremality for (some) infinite lattices. Version arXiv:2604.20947v1 (CC BY 4.0), distributed under the Creative Commons Attribution 4.0 International licence, as recorded in the arXiv licence metadata for this exact version and shown on the abstract page https://arxiv.org/abs/2604.20947v1. Full rights notice: licenses/arxiv-2604.20947-CC-BY-4.0.txt.\par\smallskip

\noindent\textit{Editorial scope.} Counted unit: the Proposition whose source label is \texttt{Prop: LM(L) is distributive sublattice} in the article (source0.tex lines 925--928, source label \texttt{Prop: LM(L) is distributive sublattice}), delivered together with the article's own complete original proof of it (source0.tex lines 930--953, one \texttt{proof} environment of 24 lines). No mathematics is written, repaired or re-derived: statement and proof are the article's own text line for line. Required context retained uncounted: Prop: Characterization of left modular (lines 809--816). Every definition, display and label of those lines is retained unchanged. Omitted from this package and not redistributed: 1--808, 817--924, 929--929, 954--1408. Nothing inside a retained excerpt is omitted or altered: every retained body is delivered exactly as the article's lines are written, so no omission receipt of any kind accompanies this package. Citations: the retained excerpts carry no citation command at all, so no bibliography entry of the article is transcribed and no reference is redistributed. Reference adaptation: none was needed and none was made; every reference inside the delivered unit resolves inside it, so no unresolved reference or citation is produced. Printed slips kept literal: no printed word, symbol, hypothesis or number of the counted statement or of its proof was altered. Layout and class: the article's own class is \texttt{amsart} with packages including color, amsmath, amsthm, bm, mathrsfs, amscd, tikz, ytableau, comment, mathdots, xy, graphicx, todonotes, enumerate; the wrapper is the lane's pinned \texttt{amsart} editorial wrapper at 10pt on A4, which restates the theorem-like environments the retained text needs and adds the article's own macro definitions that the retained text needs (\texttt{\textbackslash JI}, \texttt{\textbackslash LM}, \texttt{\textbackslash MI}), with extra wrapper packages (bm, bbm, dsfont, xspace, cancel, mathtools). The delivered numbering is the unit's own. Assets: no retained span contains a figure environment, a graphics-inclusion command, a PDF-page inclusion, a table or a TikZ picture; no image file is copied into this package, the package declares no asset dependency and no image file is redistributed with it. Licence: version arXiv:2604.20947v1 of this article is distributed under the Creative Commons Attribution 4.0 International licence, as recorded in the arXiv licence metadata for this exact version and shown on the abstract page https://arxiv.org/abs/2604.20947v1. A full rights notice is delivered at licenses/arxiv-2604.20947-CC-BY-4.0.txt. Characters: every retained excerpt is pure ASCII, so the delivered engine, XeLaTeX with its default Latin Modern text font, reports no missing character.
\begin{proposition}\label{Prop: Characterization of left modular}
Let $L$ be a weakly atomic and completely semidistributive lattice.
Then, for $t\in L$, the following conditions are equivalent:
\begin{enumerate}
    \item $t$ is left modular.
    \item For each $j\in{\JI^c}(L)$, we have $j\le t$ or $t\le \kappa(j)$, but not both.
\end{enumerate}
\end{proposition}

\begin{proposition}\label{Prop: LM(L) is distributive sublattice}
Let $L$ be a weakly atomic and completely semidistributive lattice. 
Then $\LM(L)$ is a complete sublattice of $L$ and it is completely distributive.
\end{proposition}

\begin{proof}
Let $S$ be a subset of $\LM(L)$, We only prove that $x:=\bigwedge S$ is left modular again.
Fix $j\in{\JI^c}(L)$ and $m:=\kappa(j)\in{\MI^c}(L)$, and assume $x\ {\not\leq}\ m$. Then any $s\in S$ satisfies $s\ {\not\leq}\ m$. By Proposition \ref{Prop: Characterization of left modular}, we have $j\leq s$ for each $s\in S$, and hence $j\leq x$. Again by Proposition \ref{Prop: Characterization of left modular}, $x$ is left modular.

It remains to prove that $\LM(L)$ is completely distributive.
Let $\mathcal{P}({\JI^c}(L))$ be the power set of ${\JI^c}(L)$, which forms a complete lattice. Define a map
\[f:\LM(L)\to \mathcal{P}({\JI\nolimits^c}(L)),\ x\mapsto f(x):=\{j\in{\JI\nolimits^c}(L)\mid j\le x\}.\]
This is an injective map since $x=\bigvee f(x)$ holds thanks to the weakly atomic property of $L$.
%_{j\in{\JI^c}(L),\ j\le x} j$ holds for each element $x$ in a weakly atomic complete lattice $L$.

We prove that $\LM(L)$ is a complete sublattice of $\mathcal{P}({\JI^c}(L))$. Then the claim follows since $\mathcal{P}({\JI^c}(L))$ is completely distributive.
Let $S$ be a subset $S$ of ${\JI^c}(L)$. Since  $j\in{\JI^c}(L)$ satisfies $j\le\bigwedge S$ if and only if $j\le s$ for each $s\in S$.
we have
\[f(\bigwedge S)=\bigwedge_{s\in S} f(s).\]
On the other hand, by the same argument, the map
\[g:\LM(L)\to P({\MI\nolimits^c}(L)),\ x\mapsto g(x):=\{m\in{\MI\nolimits^c}(L)\mid x\le m\}\]
satisfies $g(\bigvee S)=\bigwedge_{s\in S} g(s)$.
By Proposition \ref{Prop: Characterization of left modular}, $f(x)=\kappa^{-1}({\MI^c}(L)\setminus g(x))$ holds for each $x\in\LM(L)$. Thus
\begin{align*}
f(\bigvee S)&=\kappa^{-1}({\MI\nolimits^c}(L)\setminus g(\bigvee S))=\kappa^{-1}({\MI\nolimits^c}(L)\setminus \bigwedge_{s\in S}g(s))\\
&=\bigvee_{s\in S}\kappa^{-1}({\MI\nolimits^c}(L)\setminus g(s))=\bigvee_{s\in S}f(s)
\end{align*}
holds, as desired.
\end{proof}

\end{document}
